%% ====================================================================
\section{Data and Evaluation Protocol}
\label{sec:experiments}

\subsection{Datasets}
We use five networks from three physical domains, all public.
\begin{itemize}
\item \textbf{Delhi air quality:} daily PM$_{2.5}$ from 37 CPCB stations (2019--2023), from the
E-STGCN repository~\cite{panja2026estgcn}. Training uses 2019--2022 except the last 180 days, which
form the validation set; the test year is 2023 (365 days), with 60-day inputs and a 30-day horizon.
\item \textbf{Beijing air quality:} hourly PM$_{2.5}$ from the 12-station UCI Multi-Site
dataset~\cite{zhang2017beijingair} (2013--2017), aggregated to daily; 30-day inputs, 15-day horizon.
\item \textbf{Western-US weather:} daily maximum temperature from 30 NOAA GHCN-Daily stations
(2018--2022) with native coordinates; 30-day inputs, 15-day horizon.
\item \textbf{Late-installed stations:} daily maximum temperature from 458 established and 70 late NOAA GHCN-Daily
stations of the contiguous United States (training 2014--2018, evaluation 2023--2024), described in
Section~\ref{sec:late}.
\item \textbf{Freeway traffic:} speeds from the 228-sensor Caltrans PeMSD7-M
benchmark~\cite{yu2018stgcn} (5-minute, weekdays) and the 325-sensor PeMS-BAY benchmark~\cite{li2018dcrnn}
(subsampled to 30 minutes); 48-step inputs, 12-step horizon (1 hour on PeMSD7-M, 6 hours on PeMS-BAY). Both release a distance adjacency, not coordinates, so positions are reconstructed by
classical multidimensional scaling. They are not physical coordinates.
\end{itemize}
Within each dataset the split is chronological: the last year (air quality, weather) or the last 20\%
(traffic) is the evaluation period, preceded by a validation block of 180 days or 10\%.
\textbf{Missing data.} Non-positive and implausible readings are treated as missing and filled by forward fill,
then backward fill, then the station median, so evaluation targets may contain imputed values; no mask is used.
Values are standardized with training-network statistics.
\subsection{Random Sensor Splits}
\label{sec:splits}
Each experiment is repeated over independent random splits of the sensors into a training network
$\mathcal{T}$ and a held-out set $\mathcal{R}$, with a fresh model trained on $\mathcal{T}$ for each split. All
sensors of $\mathcal{R}$ are scored in the reference-restricted mode of Section~\ref{sec:architecture} over the same
evaluation windows. Windows advance by the horizon, so targets do not overlap; $U$ is the number of (window,
horizon-step) points per sensor. The unit of replication in every confidence interval is the split. Calibration and
test sensors are drawn \emph{from $\mathcal{R}$ at random inside each split}, thousands of times, so the expectation
of Theorem~\ref{thm:coverage} is estimated without retraining.
Training/held-out sensors are 21/16 (Delhi), 6/6 (Beijing), 18/12 (weather), 168/60 (PeMSD7-M) and 235/90
(PeMS-BAY), over 10 splits for the first three networks and 4 for the traffic networks; the supplement lists
the windows, points per sensor and calibration sizes.

\subsection{Methods Compared}
All methods are evaluated at the sensors of $\mathcal{R}$ on identical windows, with identical
calibration. We separate three kinds of method, because they use different information.
\begin{itemize}
\item \textbf{Trained forecasters that use the sensor's own history and its training-network
neighbors.} \emph{OmniST} (coordinates, neighbor attention, GRU); \emph{No-attention}
(coordinates and GRU); and \emph{Coordinate-free} (a GRU shared across sensors, with neither
coordinates nor attention). The last is the control that shows what the spatial components add: it is
also inductive, since it has no sensor-specific parameters.
\item \textbf{Forecasters with no training.} \emph{Persistence} (last observed value) and
\emph{window mean} (mean of the input window), both from the sensor's own history only.
\item \textbf{Concurrent-observation interpolators (nowcasting references).} \emph{Nearest
neighbor} and \emph{inverse-distance weighting} predict the sensor from the values that training-network
sensors report \emph{during the target window}, which a forecaster does not have at the forecast origin. They solve
an easier task and are a reference, not a competitor; the previous version of this work compared OmniST with them,
and with a kriging network, as if the tasks were the same.
\end{itemize}

\subsection{Calibration Rules and Metrics}
Intervals are symmetric, $\hat y\pm q$. We compare the \emph{sensor-level} rule of
Theorem~\ref{thm:coverage} (Corollary~\ref{cor:pooled}) with the \emph{pooled} rule, the usual
split-conformal quantile with the $(nU+1)$ correction. Both are computed from the errors of $n$
calibration sensors, in two settings. In the \emph{same-window} setting the calibration errors are those
of the window being evaluated; it is the setting of the theorem and a reference, not a deployable
procedure. In the \emph{forecast-time} setting only errors of earlier windows, and the new sensor's own
realized errors, may be used: static, expanding, rolling ($L=4$ windows for the daily networks, 24 for
the two traffic networks) and per-sensor ACI ($\gamma\in\{0.1,0.5\}$), evaluated from window $L$ onward. In the
same-window setting coverage is averaged over random (new sensor, calibration set) draws; at forecast time each
held-out sensor serves once as the new sensor with all others as calibration sensors. Coverage is the fraction of
the new sensor's points inside the interval; width is $2q$. Where $q=\infty$ the
interval is vacuous: coverage is 1 and the width is reported as infinite and excluded from width means.
Point accuracy is the mean absolute error over the sensors of $\mathcal{R}$. All intervals are
$t$-intervals over splits; differences between methods are paired within a split.
